% =============================================================================
% Appendix C: Full Proofs
% ORIUS/DC3S Battery-Safety Thesis
% =============================================================================
\chapter{Full Proofs}
\label{app:proofs}

This appendix is the authoritative proof surface for the defended ORIUS
theorem program, but it is not a flat headline ladder.  The active promotion
gate now distinguishes flagship theorems, flagship corollaries, flagship
lemmas, and archived-only historical text.  No draft theorem row remains in the
defended theorem register.  T9, T10, \texttt{T11\_Byzantine},
\texttt{T\_stale\_decay}, \texttt{T\_minimax}, and
\texttt{T\_sensor\_converse} are promoted only in their scoped forms, with
domain-discharge artifacts binding the mathematical contracts to Battery, AV,
and Healthcare evidence.  The battery witness remains the deepest proof chain,
so the full step-by-step derivations
below begin with the reference ladder:
\begin{enumerate}
  \item OASG existence,
  \item one-step safety preservation,
  \item the ORIUS core bound,
  \item observation necessity / no free safety,
  \item certificate validity horizon,
  \item certificate expiration bound,
  \item feasible fallback existence, and
  \item graceful degradation dominance.
\end{enumerate}
The certificate half-life result remains a corollary, and safe-budget
monotonicity is retained as a narrow runtime law rather than a broad converse.
The broader integrated theorem-like surface, including precursor results and
supporting chapter statements, is catalogued separately in
Appendix~\ref{app:integrated-theorem-surface-register}.  The active promotion
truth source is \path{reports/publication/theorem_promotion_matrix.csv}.

\paragraph{Mathematical contract vs. domain discharge.}
The proof copies below use only stated assumptions and typed obligations as
mathematical premises.  Domain discharge is evidence, not an extra hidden assumption: the discharge artifacts instantiate constants, verify
boundary-mass and mixing witnesses, and bind the proof kernels to the promoted
bounded Battery, AV, and Healthcare artifacts.  Thus T9/T10/T11 are
contract-universal, not unrestricted-global; new physical-AI domains may reuse
the proof pattern only after their adapters, safe sets, coverage witnesses, and
fallback obligations are discharged by domain-specific evidence.

%---------------------------------------------------------------------
\section{C.1\quad Battery-Domain OASG Existence}\label{sec:proof-oasg}
\begin{theorem}[restated]
Under A1, A2, A4, A11, and A12, for dropout rate $d>0$ and any optimal
arbitrage controller, $\exists\;t^*$ with $\phi_{\mathrm{obs}}(a_{t^*})=1$,
$\phi_{\mathrm{true}}(a_{t^*})=0$.
\end{theorem}
\begin{proof}
Fix $\eta>0$.  By A11, the battery witness controller can be driven to a
boundary-sensitive state under admissible demand, so there exists a step $t^*$
and an admissible released action $a_{t^*}$ such that the observed-state
forward prediction lies on the safe side of the upper SOC boundary while
remaining arbitrarily close to it:
\[
  \SOC_{\max}-\eta
  \le
  \hat{x}_{t^*} + \Delta(a_{t^*})
  \le
  \SOC_{\max}.
\]
By A2 and the standing dropout rate $d>0$, there exists an admissible
positive observation gap $\delta$ on a degraded-telemetry window.  Choose that
gap so that $\eta<\delta\le g_{\max}$ and let
\[
  x_{t^*} = \hat{x}_{t^*} + \delta.
\]
By A12, the fault process is conditionally independent of the controller
action stream given the history, so the boundary-sensitive witness step and
the admissible telemetry fault can co-occur on the same realized history.
Now choose the admissible one-step model residual $\epsilon_{t^*}=0$, which is
allowed by A1 because $0\in[-\epsilon_{\mathrm{model}},\epsilon_{\mathrm{model}}]$.

The controller evaluates the action through the observed state only, so the
observed successor satisfies
\[
  \hat{x}_{t^*+1}
  =
  \hat{x}_{t^*} + \Delta(a_{t^*})
  \le
  \SOC_{\max},
\]
hence $\phi_{\mathrm{obs}}(a_{t^*})=1$.  Using A4 for the shared one-step
dynamics, the true successor is
\[
  x_{t^*+1}
  =
  x_{t^*} + \Delta(a_{t^*}) + \epsilon_{t^*}
  =
  \hat{x}_{t^*} + \Delta(a_{t^*}) + \delta.
\]
Combining this identity with the lower side of the observed-state witness
gives
\[
  x_{t^*+1}
  \ge
  (\SOC_{\max}-\eta) + \delta
  =
  \SOC_{\max} + (\delta-\eta)
  >
  \SOC_{\max}.
\]
Therefore $\phi_{\mathrm{true}}(a_{t^*})=0$ while
$\phi_{\mathrm{obs}}(a_{t^*})=1$, which is exactly the battery witness-row
observation-action safety gap.  \qedhere
\end{proof}

\section{C.2\quad One-Step Safety Preservation}\label{sec:proof-one-step}
\begin{theorem}[restated]
Under A1--A5, A7, if $\asafe\in\Aset$ and $x_t\in\Crac$, then
$\phi(x_{t+1})=1$, where the tightened margin used to define $\Aset$
already absorbs the one-step model-error allowance.
\end{theorem}
\begin{proof}
Fix a step $t$ and assume the standing conditions A1--A5 and A7.  Because
$x_t\in\Crac$, the true state lies in the reliability-adjusted conformal tube
issued from the current observation.  In particular, membership in that tube gives
\[
  |x_t-\hat{x}_t| \le m_t,
\]
where $m_t$ is the inflation-adjusted half-width determined by A2 and A5.
This step is deterministic once the theorem is conditioned on $x_t\in\Crac$;
it does not introduce a new probability statement.  By A5, the shield uses
the absorbed tightening margin
\[
  m_t^\star = m_t + \epsilon_{\mathrm{model}}
\]
when constructing $\Aset$.  Since $\asafe\in\Aset$, the observed-state forward
prediction therefore satisfies
\[
  \SOC_{\min}+m_t^\star
  \le
  \hat{x}_t + \Delta(\asafe)
  \le
  \SOC_{\max}-m_t^\star,
\]
with $\Delta(\asafe)$ denoting the one-step battery state update implied by
the known dynamics from A4.  Now write the true next state as
\[
  x_{t+1} = x_t + \Delta(\asafe) + \epsilon_t,
\]
where the model residual $\epsilon_t$ is bounded by A1.  Substituting
$x_t = \hat{x}_t + (x_t-\hat{x}_t)$ and using $|x_t-\hat{x}_t|\le m_t$ yields
\[
  x_{t+1}
  \ge (\hat{x}_t+\Delta(\asafe)) - m_t - |\epsilon_t|
  \ge \SOC_{\min} + m_t^\star - m_t - \epsilon_{\mathrm{model}}
  =
  \SOC_{\min} + (m_t+\epsilon_{\mathrm{model}}) - m_t - \epsilon_{\mathrm{model}}
  =
  \SOC_{\min}.
\]
The same substitution on the upper side gives
\[
  x_{t+1}
  \le (\hat{x}_t+\Delta(\asafe)) + m_t + |\epsilon_t|
  \le \SOC_{\max} - m_t^\star + m_t + \epsilon_{\mathrm{model}}
  =
  \SOC_{\max} - (m_t+\epsilon_{\mathrm{model}}) + m_t + \epsilon_{\mathrm{model}}
  =
  \SOC_{\max}.
\]
Therefore
\[
  \SOC_{\min} \le x_{t+1} \le \SOC_{\max}.
\]
This is exactly the safety predicate $\phi(x_{t+1})=1$.  The argument uses
only the current-step information certified by A7 and therefore preserves
causality.  Hence any action selected from $\Aset$ while the current state is
inside $\Crac$ keeps the one-step successor in the safe set.  The theorem
stops at this one-step conditional statement; it does not itself quantify how
often $x_t\in\Crac$ occurs.
\qedhere
\end{proof}

\section{C.3\quad Battery ORIUS Core Bound}\label{sec:proof-core-bound}
\begin{theorem}[restated]
Under A1--A5, A7, and the explicit battery risk-envelope contract
$\Prob[Z_t=1\mid\mathcal{H}_t]\le\alpha(1-w_t)$, one has
$\E[V]\le\alpha\,T\,(1-\bar{w})$.
\end{theorem}
\begin{proof}
Let
\[
  V = \sum_{t=1}^{T} Z_t,
  \qquad
  Z_t = \mathbf{1}\!\bigl[\lnot \phi(x_t)\bigr],
\]
so that $V$ counts the number of true-state safety violations across the
episode.  Because each $Z_t$ is an indicator, $\E[Z_t]=\Prob[Z_t=1]$.
Applying linearity of expectation, which does not require independence across
time, gives
\[
  \E[V] = \sum_{t=1}^{T}\E[Z_t] = \sum_{t=1}^{T}\Prob[Z_t=1].
\]
The theorem-local battery contract supplies the per-step conditional envelope
\[
  \Prob[Z_t=1\mid\mathcal{H}_t] \le \alpha(1-w_t).
\]
Taking expectations of both sides gives
\[
  \Prob[Z_t=1] = \E[\Prob[Z_t=1\mid\mathcal{H}_t]] \le \alpha(1-w_t).
\]

Substituting the per-step bound into the expectation identity yields
\[
  \E[V]
  \le \sum_{t=1}^{T} \alpha(1-w_t)
  = \alpha \sum_{t=1}^{T}(1-w_t).
\]
Now expand the sum:
\[
  \sum_{t=1}^{T}(1-w_t)
  = T - \sum_{t=1}^{T} w_t
  = T\left(1 - \frac{1}{T}\sum_{t=1}^{T} w_t\right).
\]
By definition of the episode-average reliability,
\[
  \bar{w} = \frac{1}{T}\sum_{t=1}^{T} w_t.
\]
Replacing the bracketed term with $\bar{w}$ gives
\[
  \E[V] \le \alpha\,T\,(1-\bar{w}).
\]
This is the desired aggregation corollary.  The actual burden is the
per-step envelope hypothesis
$\Prob[Z_t=1\mid\mathcal{H}_t]\le\alpha(1-w_t)$; once that T3a-style
derivation is available, the displayed $\E[V]\le\alpha T(1-\bar w)$ bound
follows by linearity of expectation alone.  The derivation invokes only
finite-horizon identities and the explicit theorem-local battery risk
budget; no temporal independence assumption is introduced.  Therefore the
battery ORIUS core bound holds for any admissible horizon $T$ on the
defended replay surface whenever that per-step budget contract has been
justified.
\qedhere
\end{proof}

\section{C.4\quad Observation Necessity / No Free Safety}\label{sec:proof-no-free}
\begin{theorem}[restated]
Under A1, A2, A4, and A11, for any quality-ignorant $\pi$ and $d>0$, $\exists$
admissible fault sequence with $\TSVR(\pi)>0$.
\end{theorem}
\begin{proof}
Fix a quality-ignorant controller $\pi$.  By definition, its effective safety
margin does not depend on the runtime reliability score, so there exists a
constant margin $m^\pi\ge 0$ used across all observation qualities.  By A11,
the battery witness controller can be driven to a boundary-sensitive state
under admissible demand.  Consequently, for every $\eta>0$ there exists a
step $t^\star$ and an action released by $\pi$ such that the observed-state
forward prediction nearly saturates the controller's fixed margin on the upper
SOC boundary:
\[
  \hat{x}_{t^\star} + \Delta(\pi(\hat{x}_{t^\star}))
  \ge
  \SOC_{\max} - m^\pi - \eta.
\]

Because $d>0$, A2 admits an admissible degraded-observation gap.  Choose a
fault realization with
\[
  \delta \in (m^\pi+\eta,\; g_{\max}],
  \qquad
  x_{t^\star}=\hat{x}_{t^\star}+\delta.
\]
Now set the admissible one-step model residual to zero, which is allowed by A1.
Using A4 for the shared battery dynamics,
\[
  x_{t^\star+1}
  =
  x_{t^\star} + \Delta(\pi(\hat{x}_{t^\star})) + \epsilon_{t^\star}
  =
  \hat{x}_{t^\star} + \Delta(\pi(\hat{x}_{t^\star})) + \delta.
\]
Substituting the observed-state witness inequality gives
\[
  x_{t^\star+1}
  \ge
  \SOC_{\max} - m^\pi - \eta + \delta
  =
  \SOC_{\max} + (\delta-m^\pi-\eta)
  >
  \SOC_{\max}.
\]
The observed state therefore still appears feasible relative to the fixed
margin $m^\pi$, while the same released action produces a true-state boundary
crossing.  Hence the controller incurs at least one true-state violation on
the constructed admissible sequence, so $\TSVR(\pi)>0$.  The witness is
finite-horizon and class-scoped exactly as claimed in T4.  \qedhere
\end{proof}

\section{C.5\quad Certificate Validity Horizon}\label{sec:proof-cert-valid}
\paragraph{Proof source.} This section restates T5,
\emph{Finite-Horizon Certificate Validity}.
\paragraph{Theorem T5: Finite-Horizon Certificate Validity.}
Let \(X_t\) be the observation-consistent state set at issue time \(t\), let
\(\tilde a_{t:t+h}\) be the certified action or fallback sequence, and let
\(S\) be the domain safe-state set. Assume the forward reachable tube satisfies
\[
  \mathcal R_{t+j}(X_t,\tilde a_{t:t+j}) \subseteq S
  \qquad j=0,\ldots,h .
\]
Then every release governed by the certificate is safe for \(h\) steps unless a
certificate-invalidating event occurs.

\paragraph{Proof.}
At \(j=0\), all states consistent with the issued certificate lie in
\(\mathcal R_t(X_t)\subseteq S\). Suppose the claim holds through step \(j<h\).
By definition of the reachable tube, every state reachable after applying the
certified prefix \(\tilde a_{t:t+j+1}\) lies in
\(\mathcal R_{t+j+1}(X_t,\tilde a_{t:t+j+1})\). The containment assumption places
that set inside \(S\). Induction proves safety for all \(j=0,\ldots,h\). If a
fresh observation contradicts \(X_t\), the model, or the action sequence, the
certificate is outside the theorem hypothesis and must be renewed, repaired, or
expired by the runtime.


\section{C.6\quad Certificate Expiration Bound}\label{sec:proof-expiration}
\begin{theorem}[restated]
Under A4, A7, A9 and confidence level $\delta\in(0,1)$, the battery
certificate horizon satisfies the lower bound
\[
  \tau_t \ge
  \left\lfloor
    \frac{\delta_{\mathrm{bnd},t}^2}
    {2\sigma_d^2 \log(2/\delta)}
  \right\rfloor.
\]
\end{theorem}
\begin{proof}
For a $\sigma_d$-sub-Gaussian disturbance, the reflection-principle /
first-passage tail inequality gives
\[
  \Prob\!\bigl(|S_\Delta| \ge \delta_{\mathrm{bnd},t}\bigr)
  \le
  2\exp\!\left(
    -\frac{\delta_{\mathrm{bnd},t}^2}{2\sigma_d^2\Delta}
  \right).
\]
Demanding the right-hand side be at most $\delta$ and solving for $\Delta$
yields
\[
  \Delta \le
  \frac{\delta_{\mathrm{bnd},t}^2}
       {2\sigma_d^2 \log(2/\delta)}.
\]
Flooring gives the integer lower bound, and the runtime helper
\texttt{certificate\_expiration\_bound()} implements this confidence-aware
expression directly.
\qedhere
\end{proof}

\section{C.7\quad Certificate Half-Life Corollary}\label{sec:proof-halflife}
\begin{corollary}[restated]
The certificate half-life remains a corollary rather than part of the
flagship T1--T8 ladder.
\end{corollary}
\begin{proof}
The current battery blackout evidence supports a bounded safe-hold claim more
strongly than a finite empirically identified half-life law.  The corollary
therefore remains analytical and should be interpreted together with the more
conservative wording in Chapter~\ref{ch:cert-half-life-blackout}.  \qedhere
\end{proof}

\section{C.8\quad Feasible Fallback Existence}\label{sec:proof-fallback}
\begin{theorem}[restated]
Under A1, A4, and A8, there exists a battery piecewise fallback rule that
either safely holds on the interior or executes a boundary-aware safe-landing
recovery action; otherwise the runtime fails closed instead of asserting a
false fallback guarantee.
\end{theorem}
\begin{proof}
Zero dispatch is the constructive hold action on the interior.  If the current
SOC has enough slack against the one-step model error, the next state remains
inside the admissible interval.  If the current SOC is safe but too close to a
boundary for zero dispatch to be certified, the safe-landing controller
computes a charge or discharge action that moves the SOC toward the interior
safe zone while respecting the active battery limits.  The helper
\texttt{certify\_fallback\_existence()} now checks exactly this piecewise
surface and returns infeasible rather than over-claiming when neither case is
safe.  \qedhere
\end{proof}

\section{C.9\quad Graceful Degradation Dominance}\label{sec:proof-graceful}
\paragraph{Proof source.} This section restates T8,
\emph{Graceful Degradation Dominance with Useful Work}.
\paragraph{Theorem T8: Graceful Degradation Dominance with Useful Work.}
Fix a finite horizon and a paired admissible fault trace. Let \(P_G\) be an
ORIUS graceful-degradation policy and \(P_U\) an uncontrolled persistence policy
evaluated on that same trace. Let \(Z_t^G,Z_t^U\in\{0,1\}\) be true-state
violation indicators, and let \(W_G,W_U\ge0\) be the corresponding useful-work
totals on the horizon. For a declared operating threshold
\(\lambda\in(0,1]\), if
\[
  Z_t^G\le Z_t^U\quad \forall t
  \qquad\text{and}\qquad
  W_G\ge \lambda W_U,
\]
then \(P_G\) weakly dominates \(P_U\) in the safety--work preorder
\[
  P\preceq_{\lambda} Q
  \quad\Longleftrightarrow\quad
  \sum_t Z_t^P\le\sum_t Z_t^Q
  \;\text{ and }\;
  W_P\ge\lambda W_Q .
\]
This is a paired-trace dominance statement, not a claim that \(P_G\) is globally
optimal among all fallback policies.

\paragraph{Proof.}
Summing the pointwise inequality over the fixed horizon gives
\[
  \sum_t Z_t^G\le \sum_t Z_t^U,
\]
so the graceful policy has no larger violation count, and therefore no larger
violation rate, on the paired trace. The second hypothesis is exactly the
useful-work lower bound \(W_G\ge\lambda W_U\). These two inequalities are the
two coordinates of the stated preorder. Immediate shutdown can satisfy the
safety coordinate, but it is a T8 witness only if it also satisfies the same
useful-work threshold; otherwise it is outside the theorem's dominance claim.


\section{C.10\quad Safe-Budget Monotonicity Proposition}\label{sec:proof-budget}
\begin{proposition}[restated]
Under A1--A5, $B_t=|\Aset|\cdot\Delta t$ is monotone increasing in $w_t$.
\end{proposition}
\begin{proof}
$|\Aset|=\SOC_{\max}-\SOC_{\min}-2m_t$ and
$m_t=q_t/(w_t+\epsilon)$ is monotone decreasing in $w_t$, so $B_t$
increases with $w_t$.  \qedhere
\end{proof}

\section{C.11\quad Impossibility of Quality-Ignorant Mandatory Release (T9)}
\label{sec:proof-universal-impossibility}
\paragraph{Theorem T9: Impossibility of Quality-Ignorant Mandatory Release.}
Let \(\Pi_{\mathrm{obs}}\) be mandatory-release controllers
\(\pi:O\to A\). For observation \(o\), define
\[
  B(o)=\{x\in X:h(x)=o\},\qquad
  K(o)=\bigcap_{x\in B(o)} C(x).
\]
If \(K(o)=\varnothing\), no controller in \(\Pi_{\mathrm{obs}}\) can guarantee
true-state safety for all \(x\in B(o)\).

\paragraph{Proof.}
A mandatory observation-only controller must choose one action \(a=\pi(o)\) for
the entire ambiguity class \(B(o)\). If that action were safe for every
\(x\in B(o)\), then \(a\in\cap_{x\in B(o)}C(x)=K(o)\), contradicting
\(K(o)=\varnothing\). Therefore some state in the ambiguity class makes
\(\pi(o)\) unsafe. The claim is restricted to mandatory release; abstention,
fallback, uncertainty expansion, or denial of release are outside
\(\Pi_{\mathrm{obs}}\).


\section{C.12\quad Boundary-Indistinguishability Lower Bound (T10)}
\label{sec:proof-frontier-lower-bound}
\paragraph{Theorem T10: Boundary-Indistinguishability Lower Bound.}
Let \(x_0,x_1\) be two latent boundary states with observation laws \(P_0,P_1\).
Assume \(d_{\mathrm{TV}}(P_0,P_1)\le\epsilon\) and
\(C(x_0)\cap C(x_1)=\varnothing\). For any observation-only mandatory-release
controller \(\pi\),
\[
  \sup_{i\in\{0,1\}}\Pr_{O\sim P_i}[\pi(O)\notin C(x_i)]
  \ge \frac{1-\epsilon}{2}.
\]

\paragraph{Proof.}
Because the safe-action sets are disjoint, for every observation \(O\) the
action \(\pi(O)\) is unsafe for at least one of \(x_0,x_1\). Thus the sum of the
two state-conditional losses dominates the probability mass on which the two
experiments cannot be distinguished. Le Cam's two-point inequality bounds the
best discrimination advantage by \(d_{\mathrm{TV}}(P_0,P_1)\le\epsilon\), so the
average error over the two hypotheses is at least \((1-\epsilon)/2\). The
maximum of the two risks is at least their average, giving the stated bound.


\section{C.13\quad Typed Structural Transfer and Failure-Mode Converse (T11)}\label{sec:proof-typed-transfer}
\textbf{Theorem reference.}
The forward theorem is Theorem~\ref{thm:universality-completeness} and the
converse witness is Proposition~\ref{prop:t11-counterexamples} in
Chapter~\ref{ch:universality-completeness}.  The appendix copies the key proof
logic so the full-proofs surface remains self-contained.

\begin{proof}[Appendix proof copy]
Fix a step $t$ and condition on the causal observation history
$\mathcal{H}_t$.  Let $E_t = \{x_t \in U_t\}$ denote the event that the latent
state lies in the constructed observation-consistent set.  By the first T11
obligation,
\[
  \Prob[E_t \mid \mathcal{H}_t] \ge 1-\alpha.
\]
On $E_t$, the soundness obligation implies that every action in the tightened
safe-action set preserves safety at the next step for every state in $U_t$ and
every admissible disturbance.  If the tightened safe-action set is non-empty,
the repair-membership obligation guarantees that the released repaired action
lies inside that set and is therefore safe on $E_t$.  If the tightened
safe-action set is empty, the fallback-admissibility obligation guarantees a
safe release path from the current safe state.  In either branch,
\[
  \Prob[x_{t+1}\in\mathcal{S}\mid\mathcal{H}_t]
  \ge
  \Prob[E_t\mid\mathcal{H}_t]
  \ge
  1-\alpha.
\]
This proves the forward one-step transfer direction.  The converse is not
claimed as part of T11 itself; it is tracked separately as a structural
failure-mode witness showing that if any one of the four obligations fails,
the corresponding step in the forward proof disappears.
\end{proof}

\begin{corollary}[T10\_T11\_ObservationAmbiguitySandwich: Covered Observation-Ambiguity Optimality]
Let $B(o)=\{x:h(x)=o\}$ denote an observation ambiguity class and let
$K(o)=\bigcap_{x\in B(o)}C(x)$ denote its common safe core.  Every
mandatory-release observation-only controller satisfies
\[
  \operatorname{TSVR}(\pi)
  \ge
  \mathbb{E}_{O}\!\left[
    \min_{a\in\mathcal{A}}
    \Prob(a\notin C(X^\star)\mid O)
  \right].
\]
If ORIUS constructs a set $\mathcal{X}_t$ with
$\Prob(X_t^\star\notin\mathcal{X}_t)\le\alpha$ and releases only
$u_t\in\bigcap_{x\in\mathcal{X}_t}C(x)$, then
$\operatorname{TSVR}_{\mathrm{ORIUS}}\le\alpha$; under exact coverage, the
one-step true-state violation probability is zero.
\end{corollary}

\begin{proof}
Conditioning on $O=o$, an observation-only controller must select one action
for all states in $B(o)$.  Its conditional violation probability is
$\Prob(a\notin C(X^\star)\mid O=o)$, so no such controller can improve on the
minimum over $a$.  Averaging over observations gives the Bayes lower bound.
When $K(o)=\emptyset$, no action is safe for all states in the ambiguity
class, so any mandatory release has positive conditional risk if the unsafe
state has positive conditional mass.

For ORIUS, on the event $X_t^\star\in\mathcal{X}_t$, membership of $u_t$ in
every $C(x)$ for $x\in\mathcal{X}_t$ implies $u_t\in C(X_t^\star)$, hence no
true-state violation.  The remaining violation probability is bounded by the
coverage miss probability $\alpha$.
\end{proof}

\begin{lemma}[T11\_AV\_BrakeHold: AV Brake-Hold Runtime Lemma]
For a promoted ORIUS autonomous-vehicle replay row, suppose the emitted
domain runtime contract witness has source theorem T11, T11 status
\texttt{runtime\_linked}, no failed T11 obligations, a valid certificate,
and a passing true brake-hold postcondition in the runtime trace.  Then that
row is a bounded runtime instance of the forward T11 transfer surface for the
longitudinal brake-hold contract.
\end{lemma}

\begin{proof}
The witness records the four executable T11 obligations through the typed
contract summary and records the AV-specific true postcondition separately
as \texttt{domain\_postcondition\_passed}.  The stated hypotheses therefore
instantiate the coverage, soundness, repair-membership, and fallback
admissibility steps used in the T11 proof, and the runtime postcondition
checks the narrowed brake-hold safety predicate for the emitted next state.
The result is only row-local and contract-local; it does not assert complete
autonomous-driving closure.
\end{proof}

\begin{lemma}[T11\_HC\_FailSafeRelease: Healthcare Fail-Safe Release Runtime Lemma]
For a promoted ORIUS healthcare monitoring row, suppose the emitted domain
runtime contract witness has source theorem T11, T11 status
\texttt{runtime\_linked}, no failed T11 obligations, a valid certificate,
and a passing true fail-safe alert-release postcondition in the runtime
trace.  Then that row is a bounded runtime instance of the forward T11
transfer surface for the healthcare fail-safe release contract.
\end{lemma}

\begin{proof}
The witness records the typed T11 summary emitted by the universal runtime
and separately records the healthcare-specific true alert-release predicate.
If the T11 status is not \texttt{runtime\_linked}, if any obligation fails,
if the certificate is invalid, or if the postcondition fails, the witness is
not passing.  A passing witness therefore supplies exactly the forward
one-step T11 obligations plus the narrowed clinical-monitoring runtime
postcondition.  It does not assert regulated clinical deployment closure.
\end{proof}

\begin{lemma}[T6\_AV\_FallbackValidity: AV One-Step Fallback Certificate Lemma]
For a promoted ORIUS AV replay row, if the certificate records the
\texttt{single\_step\_fallback} validity scope, the released action is the
full-brake fail-safe action, the T11 runtime witness is linked, and the true
brake-hold postcondition passes, then the certificate has a theorem-backed
one-step validity horizon for that row.
\end{lemma}

\begin{proof}
The helper implementing the certificate semantics assigns
$H_t=1$, half-life one, and expiration at the next step only in the fallback
branch where the emitted action is the domain fail-safe action.  The runtime
witness separately checks T11 linkage and the true brake-hold postcondition.
Thus the certificate is valid only for the single emitted transition; no
multi-step validity is inferred from degraded observation.
\end{proof}

\begin{lemma}[T6\_HC\_FallbackValidity: Healthcare One-Step Fallback Certificate Lemma]
For a promoted ORIUS healthcare monitoring row, if the certificate records
the \texttt{single\_step\_fallback} validity scope, the released action is the
max-alert fail-safe action, the T11 runtime witness is linked, and the true
fail-safe alert-release postcondition passes, then the certificate has a
theorem-backed one-step validity horizon for that row.
\end{lemma}

\begin{proof}
The certificate-validity helper gives fallback releases a positive horizon
only for the domain fail-safe action and only for one step.  T11 linkage and
the healthcare postcondition are checked by the runtime witness.  Therefore
the row satisfies the formal validity predicate for the emitted transition,
while degraded observations are not promoted to multi-step certificates.
\end{proof}

\begin{lemma}[T\_EQ\_Battery\_RuntimeArtifactPackage: Battery Equal Artifact Discipline Package]
If the battery report directory contains runtime-denominator comparator,
ablation, negative-control, trace, utility, theorem, proof-appendix, and
reproducibility artifacts satisfying the equal-domain gate, then the battery
row satisfies the repository's equal artifact-discipline contract.
\end{lemma}

\begin{proof}
The equal-domain gate is an executable artifact obligation, not a new safety
theorem.  It checks that the battery evidence package contains a runtime
trace, runtime-native comparator rows, runtime-native ablation rows,
runtime-native negative controls, a proof-surface anchor, and a
reproducibility manifest.  It also checks the non-vacuity condition that the
ORIUS row has zero TSVR, certificate/witness closure, and useful work greater
than the degenerate safe-fallback comparator.  If any required artifact is
missing or proxy-labelled, the validator fails closed.  Therefore a passing
gate is exactly the stated artifact-discipline contract.
\end{proof}

\begin{lemma}[T\_EQ\_AV\_RuntimeArtifactPackage: AV Equal Artifact Discipline Package]
If the AV report directory contains runtime-denominator comparator,
ablation, negative-control, trace, utility, theorem, proof-appendix, and
reproducibility artifacts satisfying the equal-domain gate, then the bounded
AV replay row satisfies the repository's equal artifact-discipline contract.
\end{lemma}

\begin{proof}
The gate requires the AV T11 brake-hold runtime lemma and the AV T6
one-step fallback-validity lemma, then separately requires domain-native
runtime denominator rows for baselines, ablations, negative controls, utility,
and reproducibility.  The ORIUS row must have zero TSVR, valid runtime
certificates, passing runtime witnesses, and useful work greater than the
always-brake comparator.  The conclusion is limited to artifact discipline
for the promoted replay contract; it does not assert full autonomous-driving
closure.
\end{proof}

\begin{lemma}[T\_EQ\_HC\_RuntimeArtifactPackage: Healthcare Equal Artifact Discipline Package]
If the healthcare report directory contains runtime-denominator comparator,
ablation, negative-control, trace, utility, theorem, proof-appendix, and
reproducibility artifacts satisfying the equal-domain gate, then the bounded
healthcare monitoring row satisfies the repository's equal
artifact-discipline contract.
\end{lemma}

\begin{proof}
The gate requires the healthcare T11 fail-safe release runtime lemma and the
healthcare T6 one-step fallback-validity lemma, then requires native runtime
denominator rows for baselines, ablations, negative controls, utility, and
reproducibility.  The ORIUS row must have zero TSVR, valid runtime
certificates, passing runtime witnesses, and useful work greater than the
always-alert comparator.  The result is an artifact-discipline guarantee for
the promoted monitoring contract only; it does not assert regulated clinical
deployment readiness.
\end{proof}

%% ====================================================================
%%  C.14  L1 --- Reliability-Monotone Inflation
%% ====================================================================
\section{L1 --- Reliability-Monotone Inflation}\label{sec:proof-L1}
\paragraph{Lemma L1: Reliability-Monotone Inflation.}
Let \(m_t=q_t/(w_t+\epsilon)\) with \(q_t>0\), \(w_t\in[0,1]\), and
\(\epsilon>0\). Then \(m_t\) is monotonically nonincreasing in \(w_t\).

\paragraph{Proof.}
\[
\frac{\partial m_t}{\partial w_t}
=-\frac{q_t}{(w_t+\epsilon)^2}<0 .
\]
Thus increasing reliability cannot increase the margin.


%% ====================================================================
%%  C.15  L2 --- Safe-Set Antitonicity
%% ====================================================================
\section{L2 --- Safe-Set Antitonicity}\label{sec:proof-L2}
\paragraph{Lemma L2: Safe-Set Antitonicity.}
If \(X_1\subseteq X_2\), then
\[
\bigcap_{x\in X_2}C(x)\subseteq \bigcap_{x\in X_1}C(x).
\]

\paragraph{Proof.}
Any action safe for every state in the larger set \(X_2\) is safe for every
state in its subset \(X_1\). Hence the larger uncertainty set has a smaller or
equal common safe-action set.


%% ====================================================================
%%  C.16  L3 --- Intervention Threshold
%% ====================================================================
\section{L3 --- Intervention Threshold}\label{sec:proof-L3}
\paragraph{Lemma L3: Intervention Threshold.}
If a candidate action \(a_t\) lies outside
\(\cap_{x\in X_t}C(x)\), any certified ORIUS release must repair, fallback, or
deny release.

\paragraph{Proof.}
Membership in \(\cap_{x\in X_t}C(x)\) is exactly the condition for the action to
be safe for every state covered by the runtime uncertainty set. If the action is
outside that set, releasing it would violate the covered-release contract.
Therefore a certified runtime must choose a repair, fallback, or denial path.


%% ====================================================================
%%  C.17  L4 --- Observation-Ambiguity Safety Sandwich
%% ====================================================================
\section{L4 --- Observation-Ambiguity Safety Sandwich}\label{sec:proof-L4}
\paragraph{Lemma L4: Observation-Ambiguity Safety Sandwich.}
Observation-only mandatory release inherits the lower-side ambiguity risks in
T9 and T10. ORIUS covered release has upper-side risk bounded by its coverage
miss probability under the T2/T3 covered-release contract.

\paragraph{Proof.}
The lower side follows by applying T9 to empty safe cores and T10 to two-state
boundary ambiguity classes. For the upper side, ORIUS releases only actions
safe over the covered state set. Thus violation can occur only when the true
state is outside that covered set or an explicit runtime assumption is broken.
Under the T2/T3 contract this event has probability no larger than the
declared coverage miss probability. Combining the lower and upper statements
gives the sandwich.


%% ====================================================================
%%  C.18  T11\_Byzantine --- Byzantine-Robust Reliability Aggregation
%% ====================================================================
\section{T11\_Byzantine --- Byzantine-Robust Reliability Aggregation}
\label{sec:proof-T11-byz}
\paragraph{Theorem \texttt{T11\_Byzantine}: Byzantine-Robust Reliability Aggregation.}
Suppose \(n\) telemetry channels emit scores in \([0,1]\), at most
\(b<n/2\) are Byzantine, and \(\hat w_t\) is the \(b\)-trimmed mean after
sorting all reported scores and removing the \(b\) smallest and \(b\) largest
values. If the honest scores lie in an interval of radius \(\rho\) around
\(w_t^H\), then
\[
  |\hat w_t-w_t^H|\le \rho .
\]

\paragraph{Proof.}
Trimming removes the \(b\) smallest and \(b\) largest reported values. Since at
most \(b\) reports are Byzantine, no value below all honest scores can remain
after the lower trim: there are at most \(b\) such values, so they are among the
\(b\) smallest reports. Similarly, no value above all honest scores can remain
after the upper trim. All honest scores lie in
\([w_t^H-\rho,w_t^H+\rho]\) by hypothesis. Hence every retained value that
contributes to the trimmed mean lies in
\([w_t^H-\rho,w_t^H+\rho]\). Their arithmetic mean lies in the same interval,
which proves the bound. The theorem makes no claim for \(b\ge n/2\) or for
honest scores outside the stated interval.


%% ====================================================================
%%  C.19  T\_stale\_decay --- Stale-Hold Uncertainty Growth
%% ====================================================================
\section{T\_stale\_decay --- Stale-Hold Uncertainty Growth}
\label{sec:proof-T-stale}
\paragraph{Theorem \texttt{T\_stale\_decay}: Stale-Hold Uncertainty Growth.}
Assume no fresh observation arrives for \(s\) steps and the latent state obeys
\(\|x^\star_{t+k}-x^\star_{t+k-1}\|\le L\) for \(k=1,\ldots,s\). If the
uncertainty radius at time \(t\) is \(r_t\), then the conservative stale-hold
radius is \(r_{t+s}=r_t+Ls\).

\paragraph{Proof.}
By the triangle inequality,
\[
\|x^\star_{t+s}-x^\star_t\|
\le \sum_{k=1}^s \|x^\star_{t+k}-x^\star_{t+k-1}\|
\le Ls .
\]
Combining this drift radius with the existing uncertainty radius \(r_t\) gives
\(r_t+Ls\). Therefore any certified release at \(t+s\) must be safe over the
expanded set. This is a bounded-drift containment theorem, not a claim that a
particular exponential decay schedule is physically derived.


%% ====================================================================
%%  C.20  T\_minimax --- Finite Ambiguity-Class Minimax Lower Bound
%% ====================================================================
\section{T\_minimax --- Finite Ambiguity-Class Minimax Lower Bound}
\label{sec:proof-T-minimax}
\paragraph{Theorem Tminimax: Finite Ambiguity-Class Minimax Lower Bound.}
Let \(\mathcal P_\epsilon\) be the class of two-state boundary problems with
\(d_{\mathrm{TV}}(P_0,P_1)\le\epsilon\) and
\(C(x_0)\cap C(x_1)=\varnothing\). For observation-only mandatory-release
controllers,
\[
  \inf_{\pi\in\Pi_{\mathrm{obs}}}\sup_{P\in\mathcal P_\epsilon}
  \Pr_P[\pi(O)\notin C(X)] \ge \frac{1-\epsilon}{2}.
\]

\paragraph{Proof.}
Fix any \(\pi\). Applying Theorem T10 to the corresponding two-state problem
gives a member of \(\mathcal P_\epsilon\) whose risk is at least
\((1-\epsilon)/2\). Taking the supremum over \(\mathcal P_\epsilon\) preserves
this lower bound for the fixed \(\pi\), and taking the infimum over \(\pi\)
preserves it for the class. This is a finite ambiguity-class lower bound; it is
not a global minimax optimality theorem.


%% ====================================================================
%%  C.21  T\_sensor\_converse --- Sensor Necessity Under Adapter Semantics
%% ====================================================================
\section{T\_sensor\_converse --- Sensor Necessity Under Adapter Semantics}
\label{sec:proof-T-sensor}
\paragraph{Theorem \texttt{T\_sensor\_converse}: Sensor Necessity Under Adapter Semantics.}
Let \(x^j\) be a latent coordinate on which the safe-action map depends. Suppose
the adapter observation map is invariant to that coordinate: whenever two states
differ only in \(x^j\), they have the same observation. If there exist
\(x_0,x_1\) differing only in \(x^j\) such that
\[
  h(x_0)=h(x_1)=o,
  \qquad
  C(x_0)\cap C(x_1)=\varnothing,
\]
then no observation-only mandatory-release policy can certify safety over the
ambiguity class containing \(x_0,x_1\) without additional sensing, uncertainty
expansion, fallback, or denial.

\paragraph{Proof.}
By the adapter-invariance hypothesis, \(x_0\) and \(x_1\) produce the same
observation \(o\), so both states lie in the nonempty ambiguity class
\(B(o)\). Because \(C(x_0)\cap C(x_1)=\varnothing\), the common safe core over
that class is empty. T9 then rules out certification by any total
observation-only mandatory-release policy on \(B(o)\). Additional sensing,
uncertainty expansion, fallback, or denial are the escape routes because they
change the observation map or leave the mandatory-release policy class.


%% ====================================================================
%%  C.22  T\_trajectory\_PAC --- Finite-Horizon PAC Release Certificate
%% ====================================================================
\section{T\_trajectory\_PAC --- Finite-Horizon PAC Release Certificate}
\label{sec:proof-T-trajectory-PAC}
\paragraph{Theorem T\_trajectory\_PAC: Finite-Horizon PAC Release Certificate.}
Let \(Z_t\) be the true-state violation event at step \(t\). Suppose
\[
  \Pr(Z_t=1)\le \epsilon_t,\qquad t=1,\ldots,T .
\]
Then
\[
  \Pr\left(\sum_{t=1}^{T}Z_t\ge 1\right)
  \le \sum_{t=1}^{T}\epsilon_t .
\]
Therefore, if \(\sum_{t=1}^{T}\epsilon_t\le\delta\), the trajectory is
violation-free with probability at least \(1-\delta\).

\paragraph{Proof.}
The event \(\{\sum_{t=1}^{T}Z_t\ge 1\}\) is the union of the per-step violation
events \(\cup_{t=1}^{T}\{Z_t=1\}\). By the union bound,
\[
\Pr\left(\cup_{t=1}^{T}\{Z_t=1\}\right)
\le \sum_{t=1}^{T}\Pr(Z_t=1)
\le \sum_{t=1}^{T}\epsilon_t .
\]
If the sum is at most \(\delta\), then the complementary no-violation event has
probability at least \(1-\delta\). This defended surface is exactly the
finite-horizon Bonferroni certificate; it does not claim a Ville or
supermartingale strengthening.

